\subsection{Quantum information\label{QI}}

Let us collect a few facts about quantum pure and mixed states that are used
in the paper. We assume basic familiarity with the formalism of bras, kets,
density matrices, etc.; see Nielsen and Chuang~\cite{nc} for a good overview.

Given two mixed states $\rho$ and $\sigma$, their \emph{trace distance} is
defined as 
\[
D\left(  \rho,\sigma\right)  :=\frac{1}{2}\sum_{i=1}^{N}\left\vert
\lambda_{i}\right\vert\,,
\]
where $\lambda_{1},\ldots,\lambda_{N}$ are the
eigenvalues of $\rho-\sigma$. Trace distance is a metric and satisfies
$0\leq D\left(  \rho,\sigma\right)  \leq1$. Also, the \emph{fidelity}
$0\leq F\left(  \rho,\sigma\right)  \leq1$ is defined, in this paper, as the
maximum of $\left\vert \left\langle \psi|\varphi\right\rangle \right\vert
$ over all purifications $\left\vert \psi\right\rangle $ of $\rho$ and
$\left\vert \varphi\right\rangle $ of $\sigma$.\footnote{Some authors instead
define \textquotedblleft fidelity\textquotedblright\ as the maximum of
$\left\vert \left\langle \psi|\varphi\right\rangle \right\vert ^{2}$.} \ By
extension, given a subspace $S$, we let $F\left(  \rho,S\right)  $ be the
maximum of $\left\vert \left\langle \psi|\varphi\right\rangle \right\vert
$ over all purifications $\left\vert \psi\right\rangle $ of $\rho$ and all
unit vectors $\left\vert \varphi\right\rangle \in S$. Trace distance and
fidelity are related as follows~\cite{nc}:

\begin{proposition}
\label{ineq}For all mixed states $\rho,\sigma$,%
\[
D\left(  \rho,\sigma\right)  \leq\sqrt{1-F\left(  \rho,\sigma\right)  ^{2}}\,,
\]
with equality if $\rho$ or $\sigma$ is pure.
\end{proposition}

While fidelity is \emph{not} a metric, it does satisfy the following
inequality, which will be helpful in \expref{Section}{COR}.

\begin{lemma}
[\textquotedblleft Triangle Inequality\textquotedblright\ for Fidelity]%
\label{triangle}Suppose $\left\langle \psi\right\vert \rho\left\vert
\psi\right\rangle \geq1-\varepsilon$ and $\left\langle \varphi\right\vert
\sigma\left\vert \varphi\right\rangle \geq1-\varepsilon$. Then $F\left(
\rho,\sigma\right)  \leq\left\vert \left\langle \psi|\varphi\right\rangle
\right\vert +2\varepsilon^{1/4}$.
\end{lemma}

\begin{proof}
By \expref{Proposition}{ineq},%
\[
D\left(  \rho,\left\vert \psi\right\rangle \right)  \leq\sqrt{1-\left\langle
\psi\right\vert \rho\left\vert \psi\right\rangle }\leq\sqrt{\varepsilon}\,,
\]
and likewise $D\left(  \sigma,\left\vert \varphi\right\rangle \right)
\leq\sqrt{\varepsilon}$. Thus, since trace distance satisfies the triangle
inequality,%
\begin{align*}
D\left(  \rho,\sigma\right)   &  \geq D\left(  \left\vert \psi\right\rangle
,\left\vert \varphi\right\rangle \right)  -D\left(  \rho,\left\vert
\psi\right\rangle \right)  -D\left(  \sigma,\left\vert \varphi\right\rangle
\right)  \\
&  \geq\sqrt{1-\left\vert \left\langle \psi|\varphi\right\rangle \right\vert
^{2}}-2\sqrt{\varepsilon}\,.
\end{align*}
Then%
\begin{align*}
F\left(  \rho,\sigma\right)   &  \leq\sqrt{1-D\left(  \rho,\sigma\right)
^{2}}\\
&  \leq\sqrt{1-\left(  \sqrt{1-\left\vert \left\langle \psi|\varphi
\right\rangle \right\vert ^{2}}-2\sqrt{\varepsilon}\right)  ^{2}}\\
&  \leq\sqrt{\left\vert \left\langle \psi|\varphi\right\rangle \right\vert
^{2}+4\sqrt{\varepsilon}}\\[1ex]
&  \leq\left\vert \left\langle \psi|\varphi\right\rangle \right\vert
+2\varepsilon^{1/4}\,.
\end{align*}

\end{proof}
